\chapter{Every file, and why}
\label{ch:files}

You cannot hold 174 PHP files in your head, and you do not need to. What you
need is a reliable answer to one question: \emph{if I want to change X, which
folder do I open?} This chapter gives you that.

\section{The top level}

\begin{lstlisting}
D:\IncidentFlow\
|
+-- api/                  the Laravel backend  (all the rules live here)
+-- web/                  the React frontend   (what you look at)
+-- realtime/             the Node service     (live updates)
+-- infra/nginx/          the front door's configuration
+-- scripts/              operational scripts (database backup)
+-- deploy/               deployment runbooks
+-- docs/                 this book, and the screenshots in it
|
+-- docker-compose.yml        development: the nine containers
+-- docker-compose.prod.yml   production overlay (no bind mounts, no open ports)
+-- docker-compose.tls.yml    adds HTTPS for a single-server deployment
+-- Makefile                  short names for long commands
+-- README.md                 the short version of this book
\end{lstlisting}

\note{Three top-level folders --- \fpath{api}, \fpath{web}, \fpath{realtime} ---
are the three programs \emph{we} wrote. Everything else is either configuration
that wires them together or documentation. If you remember only one thing from
this chapter, remember those three.}

\section{Inside api/ --- where the rules live}

This is the largest and most important part of the system. Laravel puts almost
everything under \fpath{app/}.

\begin{center}
\begin{tabular}{@{}>{\raggedright\arraybackslash}p{4.3cm}rp{8.2cm}@{}}
\toprule
\textbf{Folder} & \textbf{Files} & \textbf{What is in it} \\
\midrule
\fpath{app/Models} & 14 & One class per database table. \fpath{Incident}, \fpath{User}, \fpath{AuditLog}. \\
\fpath{app/Http/Controllers} & 19 & The entry point for each HTTP request. \\
\fpath{app/Http/Resources} & 11 & Turns a model into the JSON the API returns. \\
\fpath{app/Http/Middleware} & 6 & Runs before a request reaches a controller. \\
\fpath{app/Http/Requests} & --- & Validates incoming data before anything touches it. \\
\fpath{app/Enums} & 9 & Fixed sets of values --- statuses, severities, roles. \\
\fpath{app/Policies} & 4 & Who is allowed to do what. \\
\fpath{app/Services} & 12 & The actual business logic, grouped by subject. \\
\fpath{app/Jobs} & 1 & Work that happens in the background. \\
\fpath{app/Console/Commands} & 3 & Things you can run by hand or on a schedule. \\
\fpath{app/Exceptions} & 6 & Named failures, and the one place errors are formatted. \\
\bottomrule
\end{tabular}
\end{center}

And outside \fpath{app/}:

\begin{center}
\begin{tabular}{@{}>{\raggedright\arraybackslash}p{4.3cm}p{9.4cm}@{}}
\toprule
\textbf{Folder} & \textbf{What is in it} \\
\midrule
\fpath{database/migrations} & 8 files that create the tables. Run in order, they build the schema from nothing. \\
\fpath{database/seeders} & Fills a fresh database with realistic demo data. \\
\fpath{routes/api.php} & The map from URL to controller. Start here when tracing anything. \\
\fpath{config/} & Settings, all reading from environment variables. \\
\fpath{tests/} & 13 test files, 114 tests. \\
\fpath{docker/} & The container's own files: entrypoint, PHP settings, nginx template. \\
\bottomrule
\end{tabular}
\end{center}

\subsection{One controller per action, not one per resource}

Open \fpath{app/Http/Controllers/Api/V1/} and you will find something unusual.

\begin{lstlisting}
AuthController.php                 IncidentSeverityController.php
AuditLogController.php             IncidentStatusController.php
HealthController.php               IncidentUpdateController.php
IncidentAssigneeController.php     MemberController.php
IncidentCommanderController.php    MetricsController.php
IncidentCommentController.php      NotificationController.php
IncidentController.php             OrganizationController.php
IncidentEventController.php        PostmortemController.php
IncidentExportController.php       RealtimeTicketController.php
                                   ServiceController.php
\end{lstlisting}

There are \emph{eight} controllers whose name begins with \fpath{Incident}. A
more common design would have one \fpath{IncidentController} with a dozen
methods on it: \fpath{updateStatus}, \fpath{changeSeverity}, \fpath{assign},
\fpath{comment}, \fpath{export}, and so on.

\note{Splitting them has a specific payoff. Changing an incident's
\emph{severity} and changing its \emph{status} need different permissions, emit
different timeline events, and notify different people. As methods on one class
they would share a file, and over time they would start sharing helper methods
and drifting into each other. As separate classes each has exactly one reason to
change, and the route file reads as a list of capabilities rather than a list of
URLs.}

\subsection{Where the real logic lives}

Controllers here are deliberately thin. They check the request, hand off, and
format the reply. The thinking happens in \fpath{app/Services}.

\begin{center}
\begin{tabular}{@{}>{\raggedright\arraybackslash}p{5.0cm}p{8.7cm}@{}}
\toprule
\textbf{Service folder} & \textbf{Responsible for} \\
\midrule
\fpath{Services/Incidents} & Creating and changing incidents; recording the timeline. Three files, and the most important code in the project. \\
\fpath{Services/Auth} & Issuing and verifying tokens, and the key pair they are signed with. \\
\fpath{Services/Notifications} & Deciding who hears about what, and writing the rows before anything is sent. \\
\fpath{Services/Realtime} & Publishing events to Redis so the realtime tier can forward them. \\
\fpath{Services/Metrics} & Computing time-to-acknowledge and time-to-resolve. \\
\fpath{Services/Audit} & Writing the append-only audit log. \\
\bottomrule
\end{tabular}
\end{center}

\warnbox{If you are ever lost in this codebase, the fastest route is:
\fpath{routes/api.php} tells you which controller handles a URL; the controller
tells you which service does the work; the service is where the answer is. Three
hops, every time. Resist the urge to grep --- following the path teaches you the
structure, and grep does not.}

\section{Inside web/ --- what you look at}

Thirty-six TypeScript files, in a conventional React layout.

\begin{center}
\begin{tabular}{@{}>{\raggedright\arraybackslash}p{4.3cm}p{9.4cm}@{}}
\toprule
\textbf{Folder} & \textbf{What is in it} \\
\midrule
\fpath{src/pages} & One file per screen. Nine of them --- and Part~II gives each its own chapter. \\
\fpath{src/components} & Pieces reused across screens: the layout, the timeline, badges, dialogs. \\
\fpath{src/lib} & The API client. Every request the browser makes goes through one file. \\
\fpath{src/auth} & Holding the signed-in user, and refreshing the token before it expires. \\
\fpath{src/hooks} & Data fetching, and the live connection to the realtime service. \\
\fpath{e2e/} & Browser tests that drive the whole stack the way a person would. \\
\fpath{docker/} & The nginx config that serves the built files. \\
\bottomrule
\end{tabular}
\end{center}

\note{\fpath{src/lib/api-client.ts} is worth finding early. \emph{Every} call
the frontend makes passes through it, which is why the token refresh, the error
shape, and the base URL are each defined exactly once. When you later read that
the API returns one consistent error envelope, this is the file that depends on
it being true.}

\section{Inside realtime/ --- the live connection}

The smallest of the three: twelve files, doing one thing --- holding connections
open and forwarding events onto them.

\begin{center}
\begin{tabular}{@{}>{\raggedright\arraybackslash}p{4.3cm}p{9.4cm}@{}}
\toprule
\textbf{File or folder} & \textbf{What it does} \\
\midrule
\fpath{src/index.ts} & Starts the server, and shuts it down carefully. \\
\fpath{src/app.ts} & The routes: \fpath{/healthz}, \fpath{/readyz}, \fpath{/stream}, \fpath{/metrics}. \\
\fpath{src/config.ts} & Reads and validates environment variables at boot. \\
\fpath{src/} (rest) & Subscribing to Redis, verifying tokens, tracking connections. \\
\bottomrule
\end{tabular}
\end{center}

\warnbox{The realtime service can \emph{verify} tokens but cannot \emph{create}
them. It is given only the public half of the signing key. This is deliberate:
the realtime tier holds thousands of open connections and is the most exposed
part of the system, so if it were ever compromised the attacker would gain the
ability to read tokens --- not to mint new ones. You will see the mechanism in
Part~II.}

\section{Everything else}

\begin{center}
\begin{tabular}{@{}>{\raggedright\arraybackslash}p{5.2cm}p{8.5cm}@{}}
\toprule
\textbf{Path} & \textbf{What it is} \\
\midrule
\fpath{infra/nginx/default.conf} & The routing table for the entire system: which URL goes to which service. Read it once and the architecture stops being abstract. \\
\fpath{infra/caddy/Caddyfile} & The HTTPS terminator, used only when deploying to a single server. \\
\fpath{scripts/backup-db.sh} & Nightly database backup. \\
\fpath{docker-compose.yml} & Defines the nine containers, their network, and their storage. \\
\fpath{Makefile} & Short aliases, so you type \fpath{make up} rather than a 90-character command. \\
\bottomrule
\end{tabular}
\end{center}

\section{A rule of thumb for finding anything}

\begin{center}
\begin{tabular}{@{}p{6.4cm}p{7.3cm}@{}}
\toprule
\textbf{You want to change\ldots} & \textbf{Open\ldots} \\
\midrule
what a screen looks like & \fpath{web/src/pages/} \\
what data a screen shows & \fpath{web/src/hooks/} \\
who is allowed to do something & \fpath{api/app/Policies/} \\
what happens when they do it & \fpath{api/app/Services/} \\
the shape of the JSON returned & \fpath{api/app/Http/Resources/} \\
a database column & \fpath{api/database/migrations/} \\
which URL reaches which code & \fpath{api/routes/api.php} \\
which service answers a URL & \fpath{infra/nginx/default.conf} \\
\bottomrule
\end{tabular}
\end{center}
